\scanpage{54}

l'application $f\unsure{\mathrel{\vdash\!\dashv}|\longmapsto}\rho_f$ définit une bijection de $F$ avec $ {}_3H$, enfin l'application $a\mathrel{\vdash\!\dashv}\rho_a$ fait de $A$ un rev. de degré $2$ de \deleted{\unsure{$2$}} l'ens $ {}_2H\setminus {}_2H^{+}={}_2H^{-}$ (où $H^{+}=[H,H]\simeq\unsure{A_4|\mathcal A_4|\mathfrak A_4}$) formé des él d'ordre $2$ qui ne sont pas dans le ss-gpe invariant d'ordre $4$ de $H$ ; on a $\rho_a=\rho_{a'}$ \unsure{sss|ssi} $a,a'$ égaux ou antipodiques. Les applications $s\mathrel{\vdash\!\dashv}\rho_s$ et $f\mathrel{\vdash\!\dashv}\rho_f$ transforment antipodisme en passage à l'opposé $\rho\mathrel{\vdash\!\dashv}{\unsure{\rho|g}}^{-1}$.

c) Pour tout arc $(s,\deleted{\unsure{s'}}t)$ de $\Pi$, soit $f$ l'unique face telle que $(s,\{s,t\},f)\in R^{+}$. \deleted{un rep.} Alors

\[
(\rho_s,\rho_t,\rho_{f'})\qquad(f'\ \text{l'\uline{antipodique} de }f)
\]

est un syst. de Schwartz de type $(4,4,3)$. \inserted{L'ens \deleted{de}} Ces syst. de Schwartz forment un torseur sous $H$.

\margin{$24$ d'où $24\times3=72$ syst. de Schwartz en appliquant $\unsure{\mathfrak S_3|\mathcal S_3}$}

% slanted
\margin{\uline{NB} Un ss-gpe de $H=\mathfrak S_4$ qui est d'ordre multiple de $3$ et $4$, \unsure{ou|est} \unsure{soit|est} $H^{+}=A_5$, soit $H$ (or $\rho_s,\rho_t\notin H^{+}$ !)}

Si on regarde les \inserted{$6$} transformés de $(\rho_s,\rho_t,\rho_f)$ sous $\mathfrak S_3$, se groupent en couples de $2$ ayant $=$ caractère numérique, comme

\[
(\rho_s,\rho_t,\rho_f)\text{ et }
(\underset{\substack{\shortparallel\\\rho_{t'}}}{\rho_t^{-1}},
 \underset{\substack{\shortparallel\\\rho_{s'}}}{\rho_s^{-1}},
 \underset{\substack{\shortparallel\\\rho_{f'}}}{\rho_f^{-1}}),
\]

deux tels sont conjugués sous $H$ (en effet, $(t',\{t',s'\},f')\in R^{+}$ car l'antipodisme renverse l'orientation)

d) Les syst. de Schwartz énumérés épuisent tous les cas possibles.

% archive sheet number: 53
